%! Principles of Experimental Design — Unit 1, Lesson 1.6 — Homework (Answer Key)
% PILOT SHAPE (2026-09-06): modeled on AP Statistics lesson 1.4 minus its AP Practice page.
% This is the lesson's SCORED individual practice and the LAST component in the packet.
% Students START IT IN CLASS in the last ten minutes of the period, alone, while the
% teacher circulates for the second formative read; they finish it at home. Packet pages are
% the default; a Desmos activity may replace them on a given day (decided at Close & Assign).
% THIRD CONTEXT: the notes teach on Riverbend, the Guided Practice runs on Harbor Point, and
% these items are on the Millbrook Public Library — recall is not the skill.
% TWO PAGES MAX (user decision 2026-09-06): two boxes — items 1-2 on page 1, items 3-6 and
% the spiralbox on page 2 — so no item breaks across a page.
% Item 4's rows (i) and (ii) are the deliberate CONTRAST PAIR: the same library, watched last
% summer and assigned this summer — one may claim an association, the other may say caused.
% Item 6 is the target misconception head-on (48 self-selected points vs 30 assigned points).
% THE DATA — Millbrook Public Library: 1,200 card holders; 200 volunteers -> 100 program /
%   100 waiting list; 65/100 = 65%, 35/100 = 35%, gap 30; pooled 100/200 = 50%, 0.50 x 1200 =
%   600; last summer watched: 80% vs 32%, gap 48; block on age 120 adults -> 60/60, 80
%   children -> 40/40. All verified in Python.
% PAGE LOCKSTEP: the key is generated from this file; every work block is byte-identical.
\documentclass[12pt]{article}
\usepackage{saar-article}
\usepackage{saar-key}   % pulls in -boxes; do NOT also load -boxes
\newcommand{\TallMath}[1]{$\displaystyle #1\rule[-1.4em]{0pt}{3.2em}$}
% Word bank strip for every fill-in cluster (user request 2026-09-06). Defined per document;
% the key inherits it and prints the same strip, so heights match.
\newcommand{\wordbank}[1]{\par\vspace{2pt}\noindent\fcolorbox{goldacc}{goldbg}{\parbox{\dimexpr\linewidth-2\fboxsep-2\fboxrule\relax}{\small\textbf{Word bank:} #1}}\par\vspace{4pt}}

\begin{document}

\pageheader{Unit 1, Lesson 1.6}{Homework --- Answer Key}
% NAMESTRIP: no \namedateperiod here — mirrors the blank. NO teachernote here —
% teacher prose lives in the lesson plan.

\vspace{0.12in}

% Authored identically in homework/ and homework_key/.
\begin{remindbox}
\small
\textbf{This is your graded homework.} It is scored out of 12 and due at the start of the
first class after two study halls. You will start it in the last minutes of the period, so ask about anything that stalls
you before you leave, then finish it on your own.
\end{remindbox}

\vspace{0.08in}

\boxguard[20]
\begin{scenariobox}[Millbrook Public Library --- This Time They Assign]{cerulean}
Last summer Millbrook \emph{watched}: of the card holders who had signed themselves up for
the Summer Reading Program, $80\%$ read at least $10$ books, against $32\%$ of those who had
not --- a $48$-point gap. This summer the library ($1{,}200$ card holders) runs an experiment.
$200$ card holders volunteer; a random number generator assigns $100$ to the program and
$100$ to a waiting list that starts in the fall. Everyone keeps the same card and borrowing
limits. In September, $65$ of the program group and $35$ of the waiting-list group had read at
least $10$ books.
\end{scenariobox}

\vspace{0.08in}

\boxguard
\begin{notesbox}{Practice}
Every answer names the group it is about --- the $100$ in the program, the $100$ waiting, or
all $200$ volunteers.
\wordbank{the 200 volunteers \;\textbullet\; the Summer Reading Program \;\textbullet\; the 100 on the waiting list \;\textbullet\; chance (the generator) \;\textbullet\; the director \;\textbullet\; 65\% \;\textbullet\; 35\% \;\textbullet\; 30}
\begin{enumerate}[leftmargin=1.6em, itemsep=9pt, topsep=3pt]

  \item Name the parts of this experiment.

        \par\vspace{4pt}
        Experimental units: \ans{the 200 volunteers} \hfill Treatment: \ans{the Summer Reading Program}
        \par\vspace{8pt}
        Control group: \ans{the 100 on the waiting list} \hfill Who decided the groups? \ans{chance (the generator)}

  \item Compute the reading rate for each group, then the gap.

        \textbf{(a)} What percent of the $100$ in the program read at least $10$ books?

        \begin{work}
          \dfrac{65}{100} &= 0.65 \\
                          &= 65\%
        \end{work}

        \textbf{(b)} What percent of the $100$ on the waiting list did?

        \begin{work}
          \dfrac{35}{100} &= 0.35 \\
                          &= 35\%
        \end{work}

        \textbf{(c)} The program group is higher by \ans{30} percentage points.

\end{enumerate}
\end{notesbox}

\boxguard[30]
\begin{notesbox}{Practice, continued}
\begin{enumerate}[leftmargin=1.6em, itemsep=5pt, topsep=2pt, start=3]

  \item Put all $200$ volunteers together. How many read at least $10$ books, and what
        percent is that? Then, if the volunteers represent the library, about how many of
        all $1{,}200$ card holders would that be?

        \begin{work}
          65 + 35 &= 100 \text{ card holders}; \quad \dfrac{100}{200} = 0.50 = 50\% \\
          0.50 \times 1200 &= 600 \text{ card holders --- an estimate}
        \end{work}

  \item The same library, two summers. Write what each study may claim about the program
        --- \emph{an association} or \emph{caused it} --- and the reason.

        \begin{center}
        \renewcommand{\arraystretch}{1.3}
        \begin{tabularx}{\linewidth}{@{} p{3.6cm} p{2.6cm} >{\raggedright\arraybackslash}X @{}}
        \toprule
        \textbf{The study} & \textbf{May claim} & \textbf{Because\dots} \\
        \midrule
        Last summer --- watched; $80\%$ vs.\ $32\%$
          & \ans{association} & \ans{the card holders chose their own group} \\
        This summer --- assigned; $65\%$ vs.\ $35\%$
          & \ans{caused it} & \ans{chance built the groups, not the readers} \\
        \bottomrule
        \end{tabularx}
        \end{center}

  \item Of the $200$ volunteers, $120$ are adults and $80$ are children. The director blocks
        on age and randomizes inside each block. Fill the table, then say why she blocked
        rather than leaving age to chance.

        \begin{center}
        \renewcommand{\arraystretch}{1.3}
        \begin{tabularx}{\linewidth}{@{} X c c c @{}}
        \toprule
        \textbf{Block} & \textbf{In block} & \textbf{To program} & \textbf{To waiting list} \\
        \midrule
        Adults   & $120$ & \ans{60} & \ans{60} \\
        Children & $80$  & \ans{40} & \ans{40} \\
        \bottomrule
        \end{tabularx}
        \end{center}
        \par\ansline{Blocking makes it certain: children and adults read differently.}

  \item Last summer's study found a $48$-point gap; this summer's experiment finds only
        $30$. A board member says the bigger gap is the better evidence. Explain why the
        \emph{smaller} gap is the stronger evidence that the program works.
        \par\ansline{Last summer's 48 points were inflated --- the sign-ups were already heavy readers.}
\ansline{This summer chance built the groups, so the 30 points belong to the program alone.}

\end{enumerate}
\end{notesbox}

\boxguard[5]   % the two-line spiralbox needs about five baselines; a larger guard pushes it to a third page
\begin{spiralbox}
\textbf{Next lesson: Experiments and Observational Studies Side by Side.} Millbrook could run
an experiment; a study of smoking and lung disease cannot assign anyone to smoke. Next class
is about \emph{when} an experiment is worth the cost --- and when it is impossible or wrong to
run one.
\end{spiralbox}

\end{document}
